% GENERATED FILE. Edit canonical lessons, then run npm run generate:books.
% canonical-source-hash: fnv1a64:b5ed741c68bbe40f
% canonical-lessons: TA-C348-etir, TA-C348-atcepi, TA-C348-vatitu, TA-C348-utanpatu, TA-C348-cilar

\chapter{For and Against}
\label{ch:ta-for-and-against}
\hlchaptermodality{348}

\begin{chapteropening}
\textbf{By the end of this chapter:} \emph{I can say that someone opposes or objects to an idea.}

Complete the last lesson of chapter 348: I can say that someone opposes or objects to an idea.
\end{chapteropening}

\section[\texorpdfstring{\ta{எதிர்} (etir)}{எதிர் (etir)}]{\ta{எதிர்} (etir) — to oppose; to resist}

\label{lesson:TA-C348-etir}



\begin{quote}
Before the new one: say the Tamil for a basis, then the Tamil for an explanation.
\end{quote}

\subsection*{You'll want to know: \ta{எதிர்}}
\textbf{\ta{எதிர்}} — \emph{etir} — \textquotedblleft{}to oppose; to resist\textquotedblright{}.

\begin{quote}
\ta{பல} \ta{மாணவர்கள்} \ta{இந்தத்} \ta{திட்டத்தை} \ta{எதிர்த்தார்கள்}.
\end{quote}

\begin{quote}
\emph{pala māṇavarkaḷ intat tiṭṭattai etirttārkaḷ.}
\end{quote}

\begin{quote}
Many students opposed this plan.
\end{quote}

Past (I …): \ta{எதிர்த்தேன்}, \emph{etirttēṉ}.

What is opposed takes -\ta{ஐ}: \ta{திட்டத்தை} \ta{எதிர்}. The same word means opposite, as in \ta{எதிரே}.

\begin{cousinweb}
\textbf{\ta{எதிர்}} (\emph{etir}) is native Dravidian (DEDR 795): the noun means \textquotedblleft{}that which is in front, opposite, adverse\textquotedblright{}, and the verb \textquotedblleft{}to oppose, to face\textquotedblright{}. Tamil already uses it in \ta{எதிரே}, \textquotedblleft{}opposite\textquotedblright{}.
\end{cousinweb}

\begin{grammarlens}[title={What you've built}]
One more word for giving reasons and opinions.
\end{grammarlens}

\subsection*{Your turn}
\begin{itemize}
\raggedright
  \item \emph{Say it:} \emph{etir}
  \item \emph{Say it:} \emph{etir}, once more
  \item \emph{Say it:} \emph{pala māṇavarkaḷ intat tiṭṭattai etirttārkaḷ.} — Many students opposed this plan.
  \item \emph{Recall:} say the Tamil for a basis, then the Tamil for an explanation, then say \emph{etir} again
\end{itemize}

\begin{tcolorbox}[breakable,colback=teal!4,colframe=teal!35!black,title={Before you move on}]
What does \ta{எதிர்} mean? (\textquotedblleft{}To oppose; to resist\textquotedblright{}.) Say it once more.
\end{tcolorbox}
\section[\texorpdfstring{\ta{ஆட்சேபி} (āṭcēpi)}{ஆட்சேபி (āṭcēpi)}]{\ta{ஆட்சேபி} (āṭcēpi) — to object}

\label{lesson:TA-C348-atcepi}



\begin{quote}
Before the new one: say the Tamil for an explanation, then the Tamil for to oppose.
\end{quote}

\subsection*{You'll want to know: \ta{ஆட்சேபி}}
\textbf{\ta{ஆட்சேபி}} — \emph{āṭcēpi} — \textquotedblleft{}to object\textquotedblright{}.

\begin{quote}
\ta{அந்த} \ta{முடிவை} \ta{யாரும்} \ta{ஆட்சேபிக்கவில்லை}.
\end{quote}

\begin{quote}
\emph{anta muṭivai yārum āṭcēpikkavillai.}
\end{quote}

\begin{quote}
Nobody objected to that decision.
\end{quote}

Past (I …): \ta{ஆட்சேபித்தேன்}, \emph{āṭcēpittēṉ}.

What is objected to takes -\ta{ஐ}: \ta{முடிவை} \ta{ஆட்சேபி}. The noun is \ta{ஆட்சேபனை}, an objection.

\begin{cousinweb}
\textbf{\ta{ஆட்சேபி}} (\emph{āṭcēpi}) is a verb built on Sanskrit \emph{ākṣepa}, \textquotedblleft{}objection\textquotedblright{} (CDIAL 1032). Tamil writes the Sanskrit kṣ as \ta{ட்ச}, as in \ta{சாட்சி} from \emph{sākṣin}.
\end{cousinweb}

\begin{grammarlens}[title={What you've built}]
One more word for giving reasons and opinions.
\end{grammarlens}

\subsection*{Your turn}
\begin{itemize}
\raggedright
  \item \emph{Say it:} \emph{āṭcēpi}
  \item \emph{Say it:} \emph{āṭcēpi}, once more
  \item \emph{Say it:} \emph{anta muṭivai yārum āṭcēpikkavillai.} — Nobody objected to that decision.
  \item \emph{Recall:} say the Tamil for an explanation, then the Tamil for to oppose, then say \emph{āṭcēpi} again
\end{itemize}

\begin{tcolorbox}[breakable,colback=teal!4,colframe=teal!35!black,title={Before you move on}]
What does \ta{ஆட்சேபி} mean? (\textquotedblleft{}To object\textquotedblright{}.) Say it once more.
\end{tcolorbox}
\section[\texorpdfstring{\ta{வாதிடு} (vātiṭu)}{வாதிடு (vātiṭu)}]{\ta{வாதிடு} (vātiṭu) — to argue}

\label{lesson:TA-C348-vatitu}



\begin{quote}
Before the new one: say the Tamil for to oppose, then the Tamil for to object.
\end{quote}

\subsection*{You'll want to know: \ta{வாதிடு}}
\textbf{\ta{வாதிடு}} — \emph{vātiṭu} — \textquotedblleft{}to argue\textquotedblright{}.

\begin{quote}
\ta{அவள்} \ta{தன்} \ta{உரிமைக்காக} \ta{வாதிட்டாள்}.
\end{quote}

\begin{quote}
\emph{avaḷ taṉ urimaikkāka vātiṭṭāḷ.}
\end{quote}

\begin{quote}
She argued for her right.
\end{quote}

Past (I …): \ta{வாதிட்டேன்}, \emph{vātiṭṭēṉ}.

With -\ta{உடன்} for the other side: \ta{அவருடன்} \ta{வாதிட்டேன்}, I argued with him.

\begin{cousinweb}
\textbf{\ta{வாதிடு}} (\emph{vātiṭu}) is a compound: \emph{vātam}, \textquotedblleft{}argument\textquotedblright{}, is borrowed from Sanskrit \emph{vāda}, \textquotedblleft{}speech\textquotedblright{} (CDIAL 11507), and native \emph{iṭu}, \textquotedblleft{}to put\textquotedblright{} (DEDR 442), makes it a verb, as in \ta{முத்தமிடு}.
\end{cousinweb}

\begin{grammarlens}[title={What you've built}]
One more word for giving reasons and opinions.
\end{grammarlens}

\subsection*{Your turn}
\begin{itemize}
\raggedright
  \item \emph{Say it:} \emph{vātiṭu}
  \item \emph{Say it:} \emph{vātiṭu}, once more
  \item \emph{Say it:} \emph{avaḷ taṉ urimaikkāka vātiṭṭāḷ.} — She argued for her right.
  \item \emph{Recall:} say the Tamil for to oppose, then the Tamil for to object, then say \emph{vātiṭu} again
\end{itemize}

\begin{tcolorbox}[breakable,colback=teal!4,colframe=teal!35!black,title={Before you move on}]
What does \ta{வாதிடு} mean? (\textquotedblleft{}To argue\textquotedblright{}.) Say it once more.
\end{tcolorbox}
\section[\texorpdfstring{\ta{உடன்பாடு} (uṭaṉpāṭu)}{உடன்பாடு (uṭaṉpāṭu)}]{\ta{உடன்பாடு} (uṭaṉpāṭu) — a consensus; agreement}

\label{lesson:TA-C348-utanpatu}



\begin{quote}
Before the new one: say the Tamil for to object, then the Tamil for to argue.
\end{quote}

\subsection*{You'll want to know: \ta{உடன்பாடு}}
\textbf{\ta{உடன்பாடு}} — \emph{uṭaṉpāṭu} — \textquotedblleft{}a consensus; agreement\textquotedblright{}.

\begin{quote}
\ta{இந்த} \ta{விஷயத்தில்} \ta{எங்களுக்குள்} \ta{உடன்பாடு} \ta{இருக்கிறது}.
\end{quote}

\begin{quote}
\emph{inta viṣayattil eṅkaḷukkuḷ uṭaṉpāṭu irukkiṟatu.}
\end{quote}

\begin{quote}
We are in agreement on this matter.
\end{quote}

\ta{எனக்கு} \ta{உடன்பாடு} \ta{இல்லை}, I do not agree: the dative experiencer again.

\begin{cousinweb}
\textbf{\ta{உடன்பாடு}} (\emph{uṭaṉpāṭu}) is native: DEDR 945 lists Tamil \emph{uṭaṉ}, \textquotedblleft{}together with\textquotedblright{}, and \emph{uṭaṉ-paṭu}, \textquotedblleft{}to agree\textquotedblright{}; \emph{pāṭu} is the noun of \emph{paṭu}, so the word means \textquotedblleft{}a being-together\textquotedblright{}.
\end{cousinweb}

\begin{grammarlens}[title={What you've built}]
One more word for giving reasons and opinions.
\end{grammarlens}

\subsection*{Your turn}
\begin{itemize}
\raggedright
  \item \emph{Say it:} \emph{uṭaṉpāṭu}
  \item \emph{Say it:} \emph{uṭaṉpāṭu}, once more
  \item \emph{Say it:} \emph{inta viṣayattil eṅkaḷukkuḷ uṭaṉpāṭu irukkiṟatu.} — We are in agreement on this matter.
  \item \emph{Recall:} say the Tamil for to object, then the Tamil for to argue, then say \emph{uṭaṉpāṭu} again
\end{itemize}

\begin{tcolorbox}[breakable,colback=teal!4,colframe=teal!35!black,title={Before you move on}]
What does \ta{உடன்பாடு} mean? (\textquotedblleft{}A consensus; agreement\textquotedblright{}.) Say it once more.
\end{tcolorbox}
\section[\texorpdfstring{\ta{சிலர்} (cilar)}{சிலர் (cilar)}]{\ta{சிலர்} (cilar) — some people}

\label{lesson:TA-C348-cilar}



\begin{quote}
Before the new one: say the Tamil for to argue, then the Tamil for a consensus.
\end{quote}

\subsection*{You'll want to know: \ta{சிலர்}}
\textbf{\ta{சிலர்}} — \emph{cilar} — \textquotedblleft{}some people\textquotedblright{}.

\begin{quote}
\ta{சிலர்} \ta{இந்த} \ta{யோசனையை} \ta{ஒத்துக்கொள்ளவில்லை}.
\end{quote}

\begin{quote}
\emph{cilar inta yōcaṉaiyai ottukkoḷḷavillai.}
\end{quote}

\begin{quote}
Some people did not accept this idea.
\end{quote}

Take -\ta{ஐ} for an object: \ta{சிலரைக்} \ta{கேட்டேன்}, I asked some people.

\begin{cousinweb}
\textbf{\ta{சிலர்}} (\emph{cilar}) is native Dravidian: DEDR 1571 lists Tamil \emph{cil}, \emph{cila}, \textquotedblleft{}some, few\textquotedblright{}, and \emph{cilar}, \textquotedblleft{}a few persons\textquotedblright{}, with the personal plural -\ta{அர்}.
\end{cousinweb}

\begin{grammarlens}[title={What you've built}]
One more word for giving reasons and opinions.
\end{grammarlens}

\subsection*{Your turn}
\begin{itemize}
\raggedright
  \item \emph{Say it:} \emph{cilar}
  \item \emph{Say it:} \emph{cilar}, once more
  \item \emph{Say it:} \emph{cilar inta yōcaṉaiyai ottukkoḷḷavillai.} — Some people did not accept this idea.
  \item \emph{Recall:} say the Tamil for to argue, then the Tamil for a consensus, then say \emph{cilar} again
\end{itemize}

\begin{tcolorbox}[breakable,colback=teal!4,colframe=teal!35!black,title={Before you move on}]
What does \ta{சிலர்} mean? (\textquotedblleft{}Some people\textquotedblright{}.) Say it once more.
\end{tcolorbox}
